% Источники: exam/materials/моделирование.pdf, разделы 42--46; GitHub HumsterProgrammer/iu9-notes/8sem/Моделирование/Моделирование_лекция-13_26-03-26.md.
\section{Модели машинного обучения.}

\textbf{Машинное обучение} — это раздел теории искусственного интеллекта, предметом которого является поиск методов решения задач посредством обучения через результаты решения задач той же серии.

\subsection*{Виды задач}

\begin{itemize}
  \item \textbf{Обучение с учителем (Supervised Learning).} Есть входные данные и соответствующие им правильные ответы (метки). Модель учится предсказывать ответы.
  \begin{itemize}
    \item \textbf{Регрессия}: предсказание числового значения. Пример: предсказание цены квартиры по её площади и району.
    \item \textbf{Классификация}: определение категории объекта. Пример: определение, является ли письмо спамом.
    \item \textbf{Ранжирование}: каждому объекту $x$ из $X$ необходимо определить приоритет выдачи его в качестве результата на исходный запрос.
  \end{itemize}

  \item \textbf{Обучение без учителя (Unsupervised Learning).} Даны только входные данные, нет меток. Алгоритм сам ищет структуру в данных.
  \begin{itemize}
    \item \textbf{Кластеризация}: объединение похожих объектов. Пример: сегментация клиентов по поведению на сайте.
    \item \textbf{Снижение размерности}: уменьшение числа признаков. Пример: визуализация данных с помощью PCA.
    \item \textbf{Анализ ассоциаций}: поиск закономерностей. Пример: «люди, купившие X, также покупают Y».
  \end{itemize}

  \item \textbf{Обучение с подкреплением.} Агент взаимодействует со средой и получает награды/штрафы за действия. Он учится выбирать оптимальную стратегию.
\end{itemize}

\subsection*{Обучение функциональным зависимостям}

\textbf{Функциональная зависимость} — каждому значению $x \in X$ соответствует только одно значение $y \in Y$, задаваемой функцией $y=f(x)$.

Цель состоит в том, чтобы найти функцию $f\colon X\to Y$, которая описывает зависимость между входными данными $x \in X$ и выходными значениями $y \in Y$. Эта задача решается на основе обучающей выборки
\[
  S=\{(x_i,y_i)\mid i=1,\ldots,k\}.
\]

\textbf{Пример:} предсказание цены квартиры. Нужно предсказать стоимость квартиры на основе её характеристик: общая площадь, количество комнат, этаж, район города, расстояние до метро, год постройки дома и др. Целевая переменная — цена квартиры. Если предполагается линейная функциональная зависимость, строится модель линейной регрессии:
\[
  y = \beta_0 + \beta_1 x_1 + \cdots + \beta_6 x_6.
\]

\subsection*{Основные этапы решения задач машинного обучения}

\begin{enumerate}
  \item Изучение предметной области.
  \item Предобработка данных, формирование признаков и множеств значений признаков.
  \item Построение предсказательной модели, выбор и реализация оптимизирующего алгоритма.
  \item Решение проблем переобучения и утечки данных.
  \item Оценка качества полученного результата.
  \item Внедрение и эксплуатация модели.
\end{enumerate}

\clearpage
